\loai{Loại 1. Tính hiệu suất theo Q1 và A}
\begin{vidu}
	Một động cơ nhiệt nhận từ nguồn nóng nhiệt lượng $5000\ J$ mỗi giây và có công suất $1,5\ kW$.
	\begin{vdc}
		Nhiệt lượng toả ra cho nguồn lạnh trong mỗi giây là
		\choice
		{$6500\ J$}
		{$1500\ J$}
		{$2500\ J$}
		{\True $3500\ J$}
%		\loigiai{
%			\begin{itemize}
%				\item[A] Sai. Cộng nhiệt lượng thay vì trừ công: $Q_L = Q_N + A$.
%				\item[B] Sai. Đây là công suất chứ không phải nhiệt lượng thải.
%				\item[C] Sai. Nhầm hiệu: $5000 - 2500 = 2500\ J$ không đúng.
%				\item[D] Đúng. $Q_L = Q_N - A = 5000 - 1500 = 3500\ J$
%			\end{itemize}
%		}

	\end{vdc}
	\dotlineEX{4}
	\begin{vdc}
		Hiệu suất của động cơ trên là
		\choice
		{$3\%$}
		{\True $30\%$}
		{$70\%$}
		{$3500\%$}
%		\loigiai{
%			\begin{itemize}
%				\item[A] Sai. $3\%$ là kết quả chia sai: $150/5000$
%				\item[B] Đúng. $H = \dfrac{A}{Q_N} = \dfrac{1500}{5000} = 0{,}3 = 30\%$
%				\item[C] Sai. Đây là phần nhiệt thải, không phải hiệu suất.
%				\item[D] Sai. Vô lý và đơn vị sai — không thể có hiệu suất > 100%.
%			\end{itemize}
%		}

	\end{vdc}
	\dotlineEX{4}
\end{vidu}

\begin{vidu}%[1P6B4-2][Loai1] 
	Một động cơ nhiệt thực hiện một công 400 J khi nhận từ nguồn nóng một nhiệt lượng 1 kJ. Hiệu suất của động cơ nhiệt là
	\choice
	{35\%}
	{25\%}
	{45\%}
	{\True 40\%}
	\loigiai{Hiệu suất của động cơ: $H=\dfrac{A}{Q_N}=40\%.$
	}
	%<MyLT3>
	\haidong{4}\end{vidu}


\loai{Loại 2. Tính hiệu xuất theo Q1 và Q2}
\begin{vidu}
	Một động cơ nhiệt làm việc sau một thời gian thì tác nhân đã nhận từ nguồn nóng nhiệt lượng $Q_1=1,5.10^{6}\ J$, truyền cho nguồn lạnh nhiệt lượng $Q_2^{\prime}=1,2.10^{6}\ J$. Hiệu suất của động cơ nhiệt này là
	\choice
	{\True $20 \%$}
	{$30 \%$}
	{$40 \%$}
	{$50 \%$}
	\loigiai{
		$A=Q_1-Q_2{}^{\prime}=1,5 . 10^{6}-1,2 . 10^{6}=0,3 . 10^{6}\ J$
		
		$H=\dfrac{A}{Q_1}=\dfrac{0,3 . 10^{6}}{1,5 . 10^{6}}=0,2=20 \%$.
	}
	\haidong{4}\end{vidu}

\begin{vidu}
	Một động cơ nhiệt nhả cho nguồn lạnh $80 \%$ nhiệt lượng mà nó thu được từ nguồn nóng. Hiệu suất của động cơ là
	\choice
	{\True $20 \%$}
	{$37 \%$}
	{$50 \%$}
	{$80 \%$}
	\loigiai{
		$H=\dfrac{A}{Q_1}=\dfrac{Q_1-Q_2}{Q_1}=1-\dfrac{Q_2}{Q_1}=1-0,8=0,2=20 \%$. 	
	}
	\haidong{5}\end{vidu}
\loai{Loại 3. Tính hiệu suất theo Q2 và A}
\begin{vidu}%[1P6B4-2][Loai1] 
	Trong một chu trình của động cơ nhiệt lí tưởng, chất khí thực hiện một công bằng $2.10^{3}$ J và truyền cho nguồn lạnh một nhiệt lượng bằng $6.10^{3}$ J. Hiệu suất của động cơ đó bằng
	\choice
	{33\%}
	{80\%}
	{65\%}
	{\True 25\%}
	\loigiai{
		Hiệu suất động cơ: $H=\dfrac{A}{A+Q_L}=25\%$.
	}
	%<MyLT3>
	\haidong{4}\end{vidu}
\loai{Loại 4. Từ hiệu suất, tính các đại lượng}
\begin{vidu}%[1P6B4-2][Loai3] 
	Hiệu suất của một động cơ nhiệt là 40\%, khi nguồn nóng cung cấp một nhiệt lượng 800 J, động cơ nhiệt thực hiện một công bằng
	\choice
	{2 kJ}
	{\True 320 J}
	{800 J}
	{480 J}
	\loigiai{
		Công động cơ thực hiện: $A=H.Q_N=320\ J.$
	}
	%<MyLT3>
	\haidong{4}\end{vidu}
\begin{vidu}%[1P6-4.2B][Loai3]  
	Một động cơ nhiệt có hiệu suất 30\%. Trong mỗi chu trình làm việc, tác nhân truyền cho nguồn lạnh một nhiệt lượng 240 J. Công mà động cơ thực hiện trong mỗi chu trình bằng
	\choice
	{72 J}
	{\True 103 J}
	{560 J}
	{800 J}
	\loigiai{
		\begin{itemize}
			\item Công động cơ thực hiện trong mỗi chu trình: $H=\dfrac{A}{Q_N}=\dfrac{A}{Q_L+A} \Leftrightarrow 0,3=\dfrac{A}{240+A}\Rightarrow A=103\ J.$
		\end{itemize}
	}
	\haidong{3}\end{vidu}
\begin{vidu}
	Động cơ nhiệt làm việc giữa hai nguồn nhiệt. Nhiệt lượng tác nhân nhận của nguồn nóng trong một chu trình là $2400\ J$. Hiệu suất của động cơ là $25 \%$. Nhiệt lượng truyền cho nguồn lạnh trong một chu trình là
	\choice
	{$1200\ J$}
	{$2400\ J$}
	{$600\ J$}
	{\True $1800\ J$}
	
	\loigiai{
		\begin{itemize}
			\item $A=H.Q_1=0,25.2400=600\ J$.
			\item $A=Q_1-Q_2 \Rightarrow Q_2=Q_1-A=2400-600=1800\ J$.
		\end{itemize}
	}
	\haidong{4}\end{vidu}